% ============================================================================
% CHAPTER 2 - LITERATURE REVIEW
% File: docs/thesis/theory/ch2_land_eligibility_constraints.tex
% Land eligibility constraints: rivers, settlements, terrain, minimum area.
% Operational values: docs/thesis/methods/group6_land_eligibility.tex
% ============================================================================

\section{Land Eligibility Constraints}
\label{sec:land_eligibility_review}

Land eligibility analysis removes the areas where utility-scale solar or
wind plants cannot or should not be built, before resource and cost are
considered. Frameworks such as GLAES and ENSPRESO combine legal,
environmental, social and physical restrictions into spatial exclusion
masks, and show that the resulting potential is sensitive to the choice
and value of each criterion \cite{ryberg2018,Ruiz2019}. Three groups of
criteria dominate that sensitivity and lack a single accepted value:
protection of watercourses, distance to settlements, and terrain. A
fourth question, the minimum area that can host a project, decides how
fragmented eligible land is counted.

\subsection{Watercourses and Riparian Protection}
\label{subsec:riparian_review}

Riparian zones regulate water quality, floods and sediment and connect
aquatic and terrestrial habitats. Most assessments protect them with a
fixed buffer around mapped rivers, which ignores that legal protection
often grows with river size. Where the law defines the protected strip
by channel width, a model must first estimate that width.

Downstream hydraulic geometry relates channel width to discharge by a
power law, $W = a\,Q^{b}$ \cite{leopold1953}. The coefficients are
empirical and depend on morphology, bank material, regional hydrology and
the definition of discharge. For global applications, Andreadis et al.\
estimated bankfull widths for the HydroSHEDS network with
$W_{\mathrm{bf}} = 7.2\,Q_{\mathrm{bf}}^{0.50}$, with an uncertainty of
about $\pm 0.02$ in the exponent, where $Q_{\mathrm{bf}}$ is a bankfull
discharge of roughly a two-year return period derived from drainage area
\cite{andreadis2013}. Validation against Landsat widths showed large
errors for braided reaches, and the authors present the result as a
first-order estimate. HydroRIVERS provides a consistent global river
network with long-term mean discharge per reach, but not bankfull
discharge or observed width \cite{lehner2013}. Using mean discharge in a
bankfull relation therefore gives a width proxy that is systematically
smaller than the bankfull width.

Legal rules differ in kind between countries. The Brazilian Forest Code
defines Permanent Preservation Areas along natural perennial and
intermittent watercourses as strips measured from the edge of the
regular channel, whose width increases in steps with channel width, from
30~m for channels narrower than 10~m to 500~m for channels wider than
600~m \cite{brasil2012forestcode}. Portuguese law defines the margin of
the public water domain by fixed widths that depend on the legal class of
the water, not on its size: 50~m for the sea and for navigable waters
under maritime and port jurisdiction, 30~m for other navigable or
floatable waters and public reservoirs, and 10~m for non-navigable waters
such as torrents and intermittent streams \cite{portugal2005lei54}. Elsewhere no
national rule may exist. The estimated channel width is a hydrological
quantity; the buffer attached to it is a legal parameter of each
jurisdiction, and a rule from one country is not a default for another.

\subsection{Settlements and Population}
\label{subsec:settlement_review}

Settlements are excluded for two reasons: the land is already occupied,
and plants impose local externalities such as noise, shadow flicker and
visual intrusion for wind, and visual and land-use conflicts for PV. The
usual representation is a settlement footprint with a technology-specific
setback. Reported setbacks for wind vary widely with national rules and
context: reviews find values from a few hundred metres to well above
1~km, with about 1~km as a typical value \cite{rediske2021}, German
studies use 800~m to meet noise limits for a 3~MW reference turbine and
test larger distances \cite{drechsler2017,lehmann2021,lehmann2023}, and
regional suitability studies use about 500~m for both technologies
\cite{watson2015,uyan2013}. These are planning choices, not physical or
universal legal limits.

Population density is sometimes used instead of a footprint, with a
threshold in inhabitants per square kilometre. The two are not
equivalent: density depends on the resolution and allocation method of
the population grid, and a dense pattern of rural households can exceed
a threshold without forming a settlement. A uniform density threshold
therefore excludes very different shares of land in countries with
different settlement structures. For India, a CEEW assessment on 5\,km pixels reports
that only 29\,\% of the onshore wind potential and 27\,\% of the solar
potential lie in areas below 250 inhabitants/km$^2$, so a threshold at
that value would leave out the remaining 71 and 73\,\% of the potential
estimated by that study \cite{mallya2024ceew}.
% verify: figures read on the CEEW publication page (Sept 2024); confirm page and table in the report PDF
The Degree of Urbanisation, adopted by
the UN Statistical Commission and implemented globally in the GHS
settlement model, avoids part of this problem by classifying grid cells
with density, contiguity and minimum population together, which
separates towns and villages from dispersed rural population
\cite{schiavina2023smod}.

\subsection{Terrain, Elevation and Minimum Project Area}
\label{subsec:terrain_review}

Slope controls earthworks, drainage, foundations and access. For
utility PV, NREL used 3\,\% as a conservative limit \cite{lopez2012},
engineering work on site grading reports design slopes of about 3 to
7\,\% \cite{tao2024}, and the 2024 United States Western Solar Plan
excludes land steeper than 10\,\% \cite{blm2024westernsolar}. For wind,
slope limits crane operation, component transport and foundations; NREL
used 20\,\% \cite{lopez2012}, 10\,\% has been argued from crane
constraints \cite{vanhaaren2011}, the second ENSPRESO release uses
18\,\% \cite{elsner2023onshore}, and a recent review of European siting
studies reports wind slope exclusions from about 7 to 30\,\%
\cite{epg2024renewland}. Values reported in degrees and in
percent must be converted before comparison: 15$^\circ$ corresponds to
26.8\,\%, and 15\,\% to 8.5$^\circ$. Slope also depends on the
resolution and kind of elevation model, so a threshold is tied to the
data on which it was defined. Global 30~m models such as Copernicus
GLO-30 are surface models that include vegetation and buildings; at that
spacing, a height difference of 1~m between neighbouring pixels is a
slope of 3.3\,\%, the order of the limits used for PV. FABDEM removes forests and
buildings from GLO-30 with a machine-learning correction and gives a
terrain model closer to the bare ground \cite{hawker2022fabdem}.

Elevation affects wind power through air density, since the power in the
wind is proportional to $\rho\,v^3$, and through logistics and cold
climate. Where air density is already part of the performance
calculation, an elevation cutoff no longer represents the physics and
remains a proxy for access and construction constraints. Siting studies
commonly exclude land above 1\,800 to 2\,000~m for wind
\cite{epg2024renewland}, while ENSPRESO applies no elevation cutoff and
relies on slope and resource \cite{elsner2023onshore}; the value is a
modelling convention.

Finally, eligible pixels are not necessarily a project. NREL removed
contiguous eligible areas smaller than 1~km$^2$ for rural utility PV
\cite{lopez2012}, and the second ENSPRESO release allocates wind capacity
to a 1~km$^2$ cell only when at least 10\,\% of it, 10~ha, is suitable,
enough for one 3~MW turbine \cite{elsner2023onshore}. A minimum contiguous area and a minimum eligible share within a cell are
different filters, and both depend on the resolution of the grid.